\section{Discussion}
\label{sec:discussion}

\subsection{AI Style as a Multi-Dimensional Subspace}

The most striking finding is the gap between linear probe accuracy (93--96\%) and \diffmean single-direction accuracy (69--81\%).
For refusal, \citet{arditi2024refusal} found that a single direction captures nearly the full discriminative signal ($\sim$95\% accuracy for both methods).
For truth, \citet{marks2024geometry} reported a similar pattern.
The 12--27 percentage point gap we observe for \aistyle suggests it is fundamentally different: a composite of multiple stylistic features rather than an atomic concept.

This interpretation is consistent with what we know about AI writing style.
``Sounding like AI'' involves formality, hedging (``It is worth noting that...''), verbosity, structural markers (numbered lists, topic sentences), and affective flattening.
Each of these may have its own direction in the residual stream, composing into a subspace that linear probes can exploit but a single direction cannot fully capture.
Sparse autoencoder analysis~\citep{cunningham2023sparse} could decompose this subspace into individual interpretable features, a promising direction for future work.

\subsection{Why Does the Base Model Encode AI Style?}

\pythia has never been instruction-tuned, yet its residual stream distinguishes human from ChatGPT text at 93.1\% probe accuracy.
This means the model learned enough about natural language style during pre-training to detect the statistical signatures of AI-generated text.
Pre-training on web data likely exposes the model to both casual human writing and more formal, structured text that shares features with AI-generated prose.
The model does not need to know what ``AI text'' is; it simply encodes stylistic dimensions that happen to separate the two classes.

This finding has implications for AI text detection: even base models without alignment training contain internal representations that could serve as detection signals, potentially more robust to paraphrasing than surface-level detectors.

\subsection{Layer Differences Reflect Architectural and Training Effects}

The dramatic difference in optimal layer---93.8\% depth for \pythia versus 25\% depth for \qwen---suggests that instruction tuning fundamentally reorganizes where stylistic information is represented.
In \pythia, the \aistyle signal emerges late, consistent with the general observation that base models build high-level features in their final layers.
In \qwen, the signal peaks early, possibly because alignment training creates an early ``style routing'' signal that shapes subsequent generation.
\citet{lee2024mechanistic} found that DPO overlays rather than replaces representations; our results are consistent with alignment adding an early-layer stylistic bias that the rest of the network conditions on.

\subsection{Limitations}

\para{Single AI source.}
\hcthree contains only ChatGPT responses.
The direction we find may capture ChatGPT-specific style rather than a model-general \aistyle signal.
Testing with text from Claude, Gemini, and other models is needed to assess generality.

\para{Content confounds.}
Despite the question-paired design of \hcthree, human and ChatGPT answers may differ in topic emphasis and information content, not just style.
The direction may partially encode content differences rather than pure style.

\para{Small domain sizes.}
Three of five domains have fewer than 20 samples (\textsc{medicine}: 18, \textsc{open\_qa}: 15, \textsc{wiki\_csai}: 14), limiting the reliability of cross-domain transfer estimates for these domains.

\para{Base model steering.}
Our causal steering experiment uses \pythia, a base model that does not naturally produce \aistyle text.
This limits the interpretability of the steering results; the text-level effects would likely be larger and more informative on an instruction-tuned model.

\para{Single token position.}
We collect activations only at the last token.
The \aistyle direction may be stronger or differently structured at other positions, particularly at content-bearing tokens earlier in the sequence.

\para{No human evaluation.}
We use probe accuracy and projection values as proxies for \aistyle rather than collecting human judgments of how ``AI-like'' generated text sounds under different steering conditions.
